\documentclass[11pt,a4paper]{article}
\usepackage[utf8]{inputenc}
\usepackage[T1]{fontenc}
\usepackage[portuguese]{babel}
\usepackage[margin=2.5cm]{geometry}
\usepackage{listings}
\usepackage{verbatim}

\lstset{basicstyle=\ttfamily\footnotesize,frame=single,breaklines=true}

\title{AC1 Trabalho Myvg}
\author{Aluno Daniel Santos l70070, Aluno Bernardo Gamula l58440}
\date{}

\begin{document}
\maketitle

\section{Objectivo}

Este trabalho foi integralmente realizado pelos alunos Daniel Santos (l70070) e Bernardo Gamula (l58440) sem recurso a ferramentas de inteligência artificial.

O programa lê o ficheiro .myvg, e interpreta os comandos \texttt{p}, \texttt{l} e \texttt{r} para criar a imagem ascii pretendida.

Linhas vazias e linhas que comecam por \texttt{\#} sao ignoradas.

\section{Rotinas}

\begin{itemize}
\item \texttt{sign}, \texttt{absv}: auxiliares para o algoritmo das linhas.
\item \texttt{point}: poe um \texttt{\#} na posicao (x,y). Se cair fora da imagem nao faz nada.
\item \texttt{line}: algoritmo de Bresenham.
\item \texttt{rect}: faz 4 chamadas ao \texttt{line}.
\item \texttt{initimg}: pede memoria ao sbrk e enche tudo de espacos.
\item \texttt{readln}, \texttt{nextc}, \texttt{peekc}, \texttt{skipsp}, \texttt{pint}: o parser, le linha a linha.
\item \texttt{readf}: abre o ficheiro, le, e chama o parser em loop.
\item \texttt{prtimg}: imprime a imagem no ecra.
\item \texttt{save}: escreve a imagem no ficheiro de saida.
\end{itemize}

\section{Notas}

\textbf{Leitura}
Ler 1 linha de cada vez.
readln le byte a byte com a ecall do read (a2=1) ate ao \textbackslash{}n, e mete a linha num buffer (lbuf).
Depois o peekc/nextc andam com um indice (lpos) por cima do lbuf, por isso o peek apenas le sem consumir.



\textbf{Sem mul/div} 
O enunciado diz que sao valorizados os trabalhos que nao usam mul/div/rem. Multiplicacao foi substituida por shifts:
\begin{itemize}
\item Mul var * 10 -> \texttt{pint}: \texttt{val*10 = (val<<3) + (val<<1)}.
\item Index array de pointers (4b): \texttt{slli x, 2}.
\item O \texttt{2*e} do Bresenham: \texttt{slli x, 1}.
\end{itemize}





\textbf{Ecalls} 
Codigo foi feito para o RARS, portanto usa-se a tabela de ecalls do programa:
\begin{verbatim}
1024: open
63:   read
64:   write
57:   close
9:    sbrk
10:   exit
4:    print_str
11:   print_char
8:    read_str
\end{verbatim}


\textbf{Nota: \textbackslash{}n no fim do nome.} 
O read\_string do RARS deixa \textbackslash{}n no fim do que o utilizador escreve,
isso faz com q o open nunca encontre o ficheiro.
Procuramos primeiro \texttt{\textbackslash{}n} no nome, depois \texttt{\textbackslash{}0}.

\section{Bresenham}

Fluxograma dado no enunciado:
\begin{verbatim}
dx = abs(x2 - x1)
dy = -abs(y2 - y1)
sx = sign(x2 - x1)
sy = sign(y2 - y1)
e  = dx + dy

loop:
    point(x1, y1)
    if x1 == x2 and y1 == y2: stop
    e2 = 2 * e
    if e2 >= dy: e += dy; x1 += sx
    if e2 <= dx: e += dx; y1 += sy
\end{verbatim}



Em asm o estado e' guardado em \texttt{s0..s9} 
(x1, y1, x2, y2, dx, dy, sx, sy, e, e2), 
com push/pop.





\section{Exemplo}

Input (\texttt{test.myvg}):
(Input: "C:/..path../test.myvg")
\begin{verbatim}
10 10
p 0 0
r 7 7 9 9
l 0 9 9 0
\end{verbatim}

Output ascii (\texttt{out.txt}):
\begin{verbatim}
#        #
        #
       #
      #
     #
    #
   #
  #     ###
 #      # #
#       ###
\end{verbatim}

\section{Observacoes}

A parte mais complexa neste trabalho foi o parser.
Le 1 linha para o buffer, vai lendo char a char com um indice, o peek e next lêm e avancam nesse indice.
O algoritmo das linhas so com inteiros funciona como esperado e nao precisa de nenhuma virgula flutuante.

\end{document}
